% Part: sets-functions-relations
% Chapter: functions
% Section: functions-relations

\documentclass[../../../include/open-logic-section]{subfiles}


\begin{document}

\olfileid[pt]{sfr}{fun}{rel}

\olsection{Funções como Relações}

\begin{explain}
Uma função que mapeia !!{element}s de~$A$ para !!{element}s de~$B$
define obviamente uma relação entre $A$ e~$B$, nomeadamente a relação
que se verifica entre $x$ e $y$ se e só se $f(x) = y$. De fato, podemos
até---se nos interessa reduzir os blocos de construção da
matemática, por exemplo---\emph{identificar} a função~$f$ com esta
relação, O.s.\@, com um conjunto de pares. Isto levanta então a questão:
que relações definem funções deste modo?
\end{explain}

\begin{defn}[Gráfico de uma função] Seja $f\colon A \to B$ uma função.
O \emph{gráfico} de~$f$ é a relação $R_f \subseteq A \times B$
definida por
\[
R_f = \Setabs{\tuple{x,y}}{f(x) = y}.
\]
\end{defn}

\begin{explain}
O gráfico de uma função fica unicamente determinado, pela extensionalidade.
Além disso, a extensionalidade (sobre conjuntos) justifica de imediato o
princípio implícito de extensionalidade para funções,
segundo o qual, se $f$ e~$g$ partilham domínio e contradomínio, então são
idênticas se coincidem em todos os valores.

Do mesmo modo, se uma relação é ``funcional'', então é o gráfico de uma função.
\end{explain}

\begin{prop}\ollabel{prop:graph-function}
Seja $R \subseteq A \times B$ tal que:
\begin{enumerate}
\item Se $Rxy$ e $Rxz$, então $y = z$; e
\item para todo $x \in A$ existe algum $y \in B$ tal que $\tuple{x,
y} \in R$.
\end{enumerate}
Então $R$ é o gráfico da função $f\colon A \to B$ definida por
$f(x) = y$ se e só se $Rxy$.
\end{prop}

\begin{proof}
Suponhamos que existe $y$ tal que $Rxy$. Se existisse outro $z \neq
y$ tal que $Rxz$, a condição sobre~$R$ seria violada. Logo, se
existe $y$ tal que $Rxy$, esse $y$ é único, e assim $f$ está
bem definida. Obviamente, $R_f = R$.
\end{proof}

\begin{explain}
Toda função $f\colon A \to B$ tem um gráfico, O.s.\@, uma relação em $A
\times B$ definida por $f(x) = y$. Por outro lado, toda relação~$R
\subseteq A \times B$ com as propriedades dadas em
\olref{prop:graph-function} é o gráfico de uma função~$f \colon A \to
B$. Por causa desta ligação estreita entre funções e os seus
gráficos, podemos pensar numa função simplesmente como o seu gráfico. Noutras
palavras, funções podem identificar-se com certas relações, O.s.\@, com
certos conjuntos de tuplos. \oliflabeldef{sfr:rel:ref:sec}{Note, contudo,
que o espírito desta ``identificação'' é como em
\olref[sfr][rel][ref]{sec}: não é uma afirmação sobre a metafísica das
funções, mas a observação de que é conveniente \emph{tratar}
funções como certos conjuntos. Uma razão pela qual isto é tão conveniente é que p}{P}odemos agora considerar fazer operações semelhantes sobre
funções às que fizemos sobre relações (ver
\olref[sfr][rel][ops]{sec}). Em particular:
\end{explain}

\begin{defn}\ollabel{defn:funimage}
Seja $f \colon A \to B$ uma função com $C\subseteq A$.

A \emph{restrição} de~$f$ a~$C$ é a
função~$\funrestrictionto{f}{C}\colon C \to B$ definida por
$(\funrestrictionto{f}{C})(x) = f(x)$ para todo $x \in C$. Noutras
palavras, $\funrestrictionto{f}{C} = \Setabs{\tuple{x, y} \in R_f}{x \in
C}$.

A \emph{aplicação} de~$f$ a~$C$ é $\funimage{f}{C} =
\Setabs{f(x)}{x \in C}$. Chamamos-lhe também a \emph{imagem} de~$C$
por~$f$.
\end{defn}

\begin{explain}
Segue-se destas definições que $\ran{f} =
\funimage{f}{\dom{f}}$, para qualquer função~$f$.
\oliflabeldef{sfr:rel:ops:sec}{Estas noções são exatamente as que se
esperaria, dadas as definições na \olref[sfr][rel][ops]{sec} e a nossa
identificação de funções com relações. Mas outras duas
operações---inversas e produtos relativos---exigem um pouco mais de
pormenor. Trataremos disso na \olref[inv]{sec} e
\olref[cmp]{sec}.}{}
\end{explain}

\end{document}
